\documentclass[11pt]{article}
\usepackage[margin=1.05in]{geometry}
\usepackage{amsmath,amssymb,amsthm,mathtools}
\usepackage{hyperref}

\newtheorem{theorem}{Theorem}[section]
\newtheorem{proposition}[theorem]{Proposition}
\newtheorem{lemma}[theorem]{Lemma}
\newtheorem{corollary}[theorem]{Corollary}
\theoremstyle{definition}
\newtheorem{remark}[theorem]{Remark}
\newtheorem{conjecture}[theorem]{Conjecture}

\newcommand{\Z}{\mathbb Z}
\newcommand{\R}{\mathbb R}
\newcommand{\PFtwo}{\mathrm{PF}_2}
\newcommand{\Bx}{\mathrm{Box}}
\DeclarePairedDelimiter{\abs}{\lvert}{\rvert}
\DeclareMathOperator{\supp}{supp}

\title{Concentricity is not an $m=2$ accident:\\
the crossed-sum identity and the tent lemma at every $m$}
\author{Clio}
\date{2 October 2026}

\begin{document}
\maketitle

\begin{abstract}
The proof of condition~(A) at $m=2$
(\texttt{2026-10-01-c2-inter-region-inequality.tex}, henceforth \textbf{[c2]}) turned on
\emph{concentricity}: every summand of the slice sum is symmetric about one centre
$\sigma_0$ independent of the summand. This paper settles whether that survives to
$m\ge3$. It does, for every $m$, and the mechanism is smaller than the one [c2] names.

\textbf{Theorem~\ref{thm:C}}: $\sum_{i}(L_i(\nu)+R_i(\nu))=S_\nu+\abs{\lambda}$ for every
$\nu$ and every $m$, so on a slice every summand is symmetric about the common centre
$(d-b)/2$. The two hypotheses that [c2] spends once each at $m=2$ are the $i=1$ and $i=2$
instances of \emph{one} identity, $\max(x,y+1)+\min(x-1,y)=x+y$, which is the
horizontal-strip condition read twice; cylindric periodicity is spent exactly once more,
globally, to make the single unmatched boundary term $R_m-R_0$ equal the constant $n$. A
four-way control confirms this attribution: the two deformations the proof consumes make
(C) fail, and the two it does not are silent.

\textbf{Theorem~\ref{thm:tent}} generalises the tent lemma. The support half-width is
$\Lambda(\nu)=\sigma_0-\sum_i\max(\nu_i,\lambda_{i-1}+1)$, \emph{separably} concave and
$1$-Lipschitz in each coordinate, and $\{\nu:f_\nu\ne0\}$ is again a \emph{box}; so the
effective index set is a box-slice, exchange-connected, and the occurring half-widths form
a gap-free set of consecutive integers.

What does \emph{not} survive is the certificate. At $m=2$ the layer profile $\beta$ was the
positive part of a concave function; at $m\ge3$ it is not
($24\,192$ of $37\,448$ single splines fail, and $162$ aggregate slices fail, smallest
$n=9$, $\mu=(0,1,6)$, $\lambda=(2,5,8)$, $b=2$). Worse, the route \emph{through} $\beta$ is
strictly stronger than the target and is false in the natural abstract class, while the
target survives there. The concavity reappears one level up: the half-width multiplicity
profile $M(r)=\#\{\nu:\Lambda(\nu)=r\}$ is the positive part of a concave function on all
$135\,422$ real slices tested. But \textbf{Theorem~\ref{thm:blind}} shows that even this
cannot close the argument: two multisets of concentric box splines with the \emph{same}
half-width profile $M=(1,1,1,1)$ sum to $(1,3,4,6,4,3,1)\notin\PFtwo$ and
$(1,3,6,10,6,3,1)\in\PFtwo$ respectively. So the two theorems proved here are provably
\emph{not} a complete reduction, and any proof of (A) at $m\ge3$ must use the width vectors
$w(\nu)$ and not only the half-widths they determine.
\end{abstract}

\section{Setting and statements}

Fix $m\ge1$, a level $n\ge m$, and cylindric shapes $\mu\subseteq\lambda$ in
$\mathcal C^{n,m}$, stored as bead positions $\mu_1<\dots<\mu_m<\mu_1+n$ with
$\mu_{i+m}=\mu_i+n$, likewise $\lambda$. Put $d=\abs{\lambda/\mu}=\sum_i(\lambda_i-\mu_i)$,
$\ell=3$, and $k(a,b)=K^c_{\lambda/\mu,(a,b,d-a-b)}$. Write $\PFtwo$ for: nonnegative,
support an interval of $\Z$, and $g(s)^2\ge g(s-1)g(s+1)$ for all $s$.

By [c1, Prop.~2.3] (\texttt{A-slice-sum-reformulation}, \texttt{proved}), valid for every
$m$: with
\begin{equation}\label{eq:setup}
  \Bx(\mu)=\prod_{i=1}^m[\mu_i,\mu_{i+1}-1],\qquad
  L_i(\nu)=\max(\nu_i,\lambda_{i-1}+1),\qquad
  R_i(\nu)=\min(\nu_{i+1}-1,\lambda_i),
\end{equation}
\begin{equation}\label{eq:slice}
  k(a,b)=\!\!\sum_{\substack{\nu\in\Bx(\mu)\\ S_\nu=\abs{\mu}+b}}\!\! f_\nu(a),
  \qquad
  f_\nu=\mathop{\text{\Large$*$}}_{i=1}^{m}\mathbf 1_{[\,L_i(\nu)-\nu_i,\ R_i(\nu)-\nu_i\,]},
  \qquad S_\nu=\textstyle\sum_i\nu_i,
\end{equation}
where throughout \emph{indices are cyclic}: $\lambda_{j+m}=\lambda_j+n$ and
$\nu_{j+m}=\nu_j+n$, so in particular $\lambda_0=\lambda_m-n$ and $\nu_{m+1}=\nu_1+n$.
Values of $\nu$ with $L_i(\nu)>R_i(\nu)$ for some $i$ contribute $f_\nu=0$. Write
\[
  w_i(\nu)=R_i(\nu)-L_i(\nu)+1,\qquad
  \Sigma_b=\{\nu\in\Bx(\mu):S_\nu=\abs{\mu}+b\},\qquad
  \Sigma_b^{\circ}=\{\nu\in\Sigma_b: f_\nu\ne0\}.
\]
Condition~(A) is the assertion that $a\mapsto k(a,b)$ is $\PFtwo$ for every $b$; it is
\texttt{proved} at $m=2$ [c2, Thm.~1.1] and \texttt{computed} (open) at $m\ge3$.

The brief for this session named one test, and it is an equivalence rather than a heuristic.
A convolution $*_i\mathbf 1_{[\alpha_i,\beta_i]}$ is symmetric about
$\bigl(\sum_i\alpha_i+\sum_i\beta_i\bigr)/2$; here $\alpha_i=L_i-\nu_i$ and
$\beta_i=R_i-\nu_i$, so $f_\nu$ is symmetric about
\begin{equation}\label{eq:centre}
  c(\nu)=\tfrac12\Bigl(\textstyle\sum_i\bigl(L_i(\nu)+R_i(\nu)\bigr)-2S_\nu\Bigr),
\end{equation}
and since $S_\nu$ is constant on a slice,
\begin{equation}\label{eq:C}
  \text{(C)}\qquad \textstyle\sum_i\bigl(L_i(\nu)+R_i(\nu)\bigr)\ \text{is independent of }
  \nu\ \text{on each slice}
\end{equation}
holds if and only if all summands on the slice share a centre.

\begin{theorem}[(C), with its value, for every $m$]\label{thm:C}
For every $m\ge1$, every cylindric $\mu\subseteq\lambda$ in $\mathcal C^{n,m}$ and every
$\nu\in\Z^m$,
\[
  \sum_{i=1}^m\bigl(L_i(\nu)+R_i(\nu)\bigr)\ =\ S_\nu+\abs{\lambda}.
\]
Consequently every $f_\nu$ with $\nu\in\Sigma_b$ is symmetric about the common centre
\[
  c\ =\ \tfrac12\bigl(\abs{\lambda}-\abs{\mu}-b\bigr)\ =\ \tfrac{d-b}{2},
\]
and $k(\cdot,b)$ is symmetric about $(d-b)/2$. Concentricity is not an $m=2$ accident.
\end{theorem}

\begin{theorem}[tent lemma at every $m$]\label{thm:tent}
Put $\sigma_0(\nu)=\tfrac12\bigl(S_\nu+\abs{\lambda}\bigr)$. For every $\nu$ with
$f_\nu\ne0$ the support of $f_\nu$ is $[\,c-\Lambda(\nu),\ c+\Lambda(\nu)\,]$ where
\begin{equation}\label{eq:Lam}
  \Lambda(\nu)\ =\ \sigma_0(\nu)-\sum_{i=1}^m\max\bigl(\nu_i,\ \lambda_{i-1}+1\bigr)
  \ =\ \tfrac12 S_\nu-\tfrac12\abs{\lambda}+n-m-\sum_{i=1}^m\bigl(\nu_i-\lambda_{i-1}-1\bigr)_+ .
\end{equation}
Moreover:
\begin{enumerate}
\item On a slice, $-\Lambda$ is a sum of $m$ convex, nondecreasing, $1$-Lipschitz functions
of the \emph{single} coordinates $\nu_1,\dots,\nu_m$ respectively; in particular $\Lambda$
is concave on $\R^m$ and $1$-Lipschitz in each coordinate.
\item $\Sigma_b^\circ=\Sigma_b\cap\Bx^\circ$ where $\Bx^\circ$ is again a box:
\begin{equation}\label{eq:effbox}
  \Bx^{\circ}=\prod_{i=1}^m\bigl[\max(\mu_i,\lambda_{i-2}+2),\ \min(\mu_{i+1}-1,\lambda_i)\bigr].
\end{equation}
\item Hence $\Sigma^\circ_b$ is connected under the exchange moves
$\nu\mapsto\nu+e_i-e_j$, along each of which $\abs{\Delta\Lambda}\le1$, so
$\Lambda(\Sigma^\circ_b)$ is a set of \emph{consecutive} integers in the coset
$\sigma_0+\Z$ --- it has no gaps --- and $\Lambda$ attains its maximum on $\Sigma^\circ_b$.
\end{enumerate}
At $m=2$ this is [c2, Lem.~2.4] with $\Lambda=L_t$.
\end{theorem}

Theorems~\ref{thm:C} and~\ref{thm:tent} are steps 1 and 2 of the $m=2$ architecture,
generalised. Section~\ref{sec:dead} shows that step~3 --- $\beta$ is the positive part of a
concave function --- is \emph{false} at $m\ge3$, and that the natural repair (prove $\beta$
log-concave and quote the tail lemma) is strictly stronger than the target and false in the
abstract class the two theorems above define. Section~\ref{sec:gap} isolates what remains.

\section{Proof of Theorem~\ref{thm:C}}\label{sec:proofC}

Everything is one scalar identity.

\begin{lemma}[the adjacency identity]\label{lem:adj}
For all $x,y\in\Z$ and all $p\in\Z$,
\[
  \max(x,\ y+p)\ +\ \min(x-p,\ y)\ =\ x+y .
\]
\end{lemma}
\begin{proof}
If $x\le y$ then $x-p\le y$ and $y+p\ge x$, so the left side is $(y+p)+(x-p)=x+y$. If
$x\ge y+p$ then $x-p\ge y$, so the left side is $x+y$. If $y<x<y+p$ then
$\max(x,y+p)=y+p$ and $x-p<y$ so $\min(x-p,y)=x-p$, again $x+y$.
\end{proof}

The identity is false as soon as the two offsets differ: $\max(x,y+p)+\min(x-q,y)$ equals
$x+y+p-q$ for $x\ll y$ and $x+y$ for $x\gg y$. The matched offset $p=q$ is therefore the
load-bearing hypothesis, and $p=q=1$ is exactly what \eqref{eq:setup} supplies.

\begin{lemma}[crossed sums at every $m$]\label{lem:crossed}
For every $i\in\Z$ and every $\nu$,
\[
  L_i(\nu)+R_{i-1}(\nu)\ =\ \nu_i+\lambda_{i-1},
\]
where $L,R$ are extended to all $i\in\Z$ by \eqref{eq:setup} and the cyclic conventions.
\end{lemma}
\begin{proof}
$L_i=\max(\nu_i,\ \lambda_{i-1}+1)$ and $R_{i-1}=\min(\nu_i-1,\ \lambda_{i-1})$. Apply
Lemma~\ref{lem:adj} with $x=\nu_i$, $y=\lambda_{i-1}$, $p=q=1$.
\end{proof}

The two clamps paired here are $(\nu_i,\nu_i-1)$ and $(\lambda_{i-1}+1,\lambda_{i-1})$: each
is a pair of \emph{consecutive integers coming from one number}. That is not a coincidence
of the formulas but the content of the two horizontal-strip conditions,
$\nu_i\le\kappa_i<\nu_{i+1}$ and $\kappa_{i-1}\le\lambda_{i-1}<\kappa_i$. So
Lemma~\ref{lem:crossed} is the horizontal-strip condition read twice, once from each side,
and it is bead-local: it knows nothing about $m$.

\begin{proof}[Proof of Theorem~\ref{thm:C}]
Pair each $L_i$ with $R_{i-1}$ for $i=1,\dots,m$. The sum $\sum_{i=1}^mR_{i-1}$ differs
from $\sum_{i=1}^mR_i$ by exactly the two unmatched terms, so
\begin{equation}\label{eq:tele}
  \sum_{i=1}^m\bigl(L_i+R_i\bigr)
  \;=\;\sum_{i=1}^m\bigl(L_i+R_{i-1}\bigr)\;+\;\bigl(R_m-R_0\bigr)
  \;=\;\sum_{i=1}^m\bigl(\nu_i+\lambda_{i-1}\bigr)\;+\;\bigl(R_m-R_0\bigr)
\end{equation}
by Lemma~\ref{lem:crossed}. Now the cyclic conventions close the cycle twice. First,
$R_0(\nu)=\min(\nu_1-1,\lambda_0)=\min(\nu_1+n-1,\lambda_m)-n=R_m(\nu)-n$, so
$R_m-R_0=n$ --- a constant. Second,
$\sum_{i=1}^m\lambda_{i-1}=\lambda_0+\lambda_1+\dots+\lambda_{m-1}
=\abs{\lambda}-\lambda_m+(\lambda_m-n)=\abs{\lambda}-n$. Substituting both
into~\eqref{eq:tele},
\[
  \sum_{i=1}^m\bigl(L_i+R_i\bigr)=S_\nu+\abs{\lambda}-n+n=S_\nu+\abs{\lambda}.
\]
For the consequence: by~\eqref{eq:centre} the centre of $f_\nu$ is
$\tfrac12\bigl(S_\nu+\abs{\lambda}-2S_\nu\bigr)=\tfrac12(\abs{\lambda}-S_\nu)$, which on
$\Sigma_b$ equals $\tfrac12(\abs{\lambda}-\abs{\mu}-b)=(d-b)/2$, independent of $\nu$. A
sum of functions symmetric about a common point is symmetric about it.
\end{proof}

\subsection{Where each hypothesis is spent}

The discipline of [c2] --- ask \emph{where} a hypothesis is spent, not what it implies ---
gives a cleaner answer here than the one [c2] itself records.

\begin{enumerate}
\item \textbf{The matched offset $p=q=1$} (equivalently: both horizontal-strip conditions)
is spent once per bead, $m$ times in all, in Lemma~\ref{lem:crossed}.
\item \textbf{Cylindric periodicity} ($\lambda_{j+m}=\lambda_j+n$ \emph{and}
$\nu_{j+m}=\nu_j+n$, with the \emph{same} $n$) is spent exactly once, globally, on the
single unmatched boundary term, to give $R_m-R_0=n$. It is spent a second time only to
evaluate the constant, $\sum_i\lambda_{i-1}=\abs{\lambda}-n$.
\end{enumerate}

[c2] describes the $m=2$ proof as spending two hypotheses, $A=D-n+1$ and $C=B+1$, once each.
That is the $m=2$ shadow of item~1: in the dictionary
$a_t=L_1,\ b_t=R_1,\ c_t=L_2,\ d_t=R_2$ (with $t=\nu_1$, $u=S_\nu$), the identity
$b_t+c_t=u+B-t$ is the $i=2$ instance of Lemma~\ref{lem:crossed} and $a_t+d_t=t+D$ is the
$i=1$ instance together with the wrap $R_0=R_m-n$. So there is one local identity used
$m$ times, plus one global closure --- and at $m=2$ the global closure happens to be
invisible, because with only two beads it hides inside the second crossed sum.

\begin{remark}[a correction to my own first reading]
On first deriving Lemma~\ref{lem:crossed} I concluded that periodicity is \emph{not}
load-bearing for (C), on the ground that $\lambda_0$ enters only the additive constant. That
is wrong, and the control in §\ref{sec:control} said so. The error was to look at
$\sum_i\lambda_{i-1}$ and forget $R_m-R_0$: periodicity is what makes the one unmatched
pair differ by a constant. Decoupling the $\lambda$-period from the $\nu$-period breaks (C)
(control P3 below), while changing the offset \emph{value} $p=q$ does not (control P1).
\end{remark}

\section{Proof of Theorem~\ref{thm:tent}}\label{sec:prooftent}

\begin{proof}
\emph{The formula \eqref{eq:Lam}.} If $f_\nu\ne0$ then every $w_i\ge1$ and $f_\nu$ is
supported on an interval of length $\sum_iw_i-m+1$, so its half-width is
$\Lambda=\tfrac12(\sum_iw_i-m)$. By Theorem~\ref{thm:C},
$\sum_iR_i=S_\nu+\abs{\lambda}-\sum_iL_i$, hence
$\sum_iw_i=\sum_iR_i-\sum_iL_i+m=S_\nu+\abs{\lambda}-2\sum_iL_i+m$ and
\[
  \Lambda(\nu)=\tfrac12\bigl(S_\nu+\abs{\lambda}\bigr)-\sum_iL_i(\nu)
  =\sigma_0(\nu)-\sum_i\max(\nu_i,\lambda_{i-1}+1),
\]
which is the first form. For the second, write
$\max(\nu_i,\lambda_{i-1}+1)=(\lambda_{i-1}+1)+(\nu_i-\lambda_{i-1}-1)_+$ and use
$\sum_i(\lambda_{i-1}+1)=\abs{\lambda}-n+m$.

\emph{(1).} $\phi_i(x):=\max(x,\lambda_{i-1}+1)$ is convex, nondecreasing and $1$-Lipschitz
on $\Z$, with increments $\phi_i(x+1)-\phi_i(x)\in\{0,1\}$, equal to $1$ exactly when
$x\ge\lambda_{i-1}+1$. On a slice $\sigma_0$ is a constant, so
$\Lambda=\text{const}-\sum_i\phi_i(\nu_i)$ is a constant minus a separable convex function.

\emph{(2).} $f_\nu\ne0$ iff $L_i(\nu)\le R_i(\nu)$ for every $i$, i.e.\ iff
$\max(\nu_i,\lambda_{i-1}+1)\le\min(\nu_{i+1}-1,\lambda_i)$, which is the conjunction of
four inequalities:
\[
  \nu_i\le\nu_{i+1}-1,\qquad \nu_i\le\lambda_i,\qquad
  \lambda_{i-1}+1\le\nu_{i+1}-1,\qquad \lambda_{i-1}+1\le\lambda_i .
\]
The first holds automatically on $\Bx(\mu)$, since $\nu_i\le\mu_{i+1}-1\le\nu_{i+1}-1$; the
fourth holds because $\lambda$ is a shape. The second and third are each a bound on a
\emph{single} coordinate; reindexing the third ($i+1\rightsquigarrow i$) it reads
$\nu_i\ge\lambda_{i-2}+2$. Intersecting with $\Bx(\mu)$ gives~\eqref{eq:effbox}.

\emph{(3).} $\Sigma^\circ_b$ is the set of integer points of a box lying on the hyperplane
$\sum_i x_i=\abs{\mu}+b$. Such a set is connected under the moves $x\mapsto x+e_i-e_j$: if
$x\ne y$ lie in it, then $\sum_i(x_i-y_i)=0$, so there are $i$ with $x_i<y_i$ and $j$ with
$x_j>y_j$, and $x+e_i-e_j$ is again in the box (each coordinate moves one step towards
$y_i$ resp.\ $y_j$, hence stays between the box bounds) and on the hyperplane, strictly
reducing $\sum_i\abs{x_i-y_i}$. Along such a move
\[
  \Delta\Lambda=-\bigl[\phi_i(x_i+1)-\phi_i(x_i)\bigr]+\bigl[\phi_j(x_j)-\phi_j(x_j-1)\bigr]
  \in\{-1,0,1\}
\]
because each bracket lies in $\{0,1\}$. A function with increments in $\{-1,0,1\}$ along the
edges of a connected graph has an image with no gaps, and the image is finite, so the
maximum is attained. All values lie in $\sigma_0+\Z$ since each $\phi_i$ is $\Z$-valued.
\end{proof}

\begin{remark}
Part~(2) is the general-$m$ replacement for the $m=2$ interval
$T=[\max(\mu_1,u-\mu_1-n+1),\ \min(\mu_2-1,u-\mu_2)]$ of [c2, eq.~(3)]: at $m=2$ the
effective index set is a box-slice in two coordinates, i.e.\ an interval in $t=\nu_1$, and
that is where [c2] got its $T_\pm$. Part~(2) explains why no extra ``$h_t\le0$'' case
analysis is needed: the degenerate $\nu$ are cut off by \emph{box} constraints, not by a
coupled condition.
\end{remark}

\section{The $m=2$ certificate is dead at $m\ge3$}\label{sec:dead}

By Theorem~\ref{thm:C}, on a slice we may write $x=a-c$ and
\[
  G(a):=k(a,b)=\gamma\bigl(\abs{x}\bigr),\qquad \gamma(y)=G(c+y),
\]
with $\gamma$ nonnegative and nonincreasing (each $f_\nu$ is a convolution of interval
indicators, hence symmetric and unimodal about $c$). Define the \emph{layer profile}
\[
  \beta(r)\ :=\ \gamma(r)-\gamma(r+1)\ \ge\ 0,
  \qquad\text{so}\qquad
  f=\sum_r\beta(r)\,\mathbf 1_{[c-r,\,c+r]},\quad \gamma(y)=\sum_{r\ge y}\beta(r).
\]
Each individual $f_\nu$ has its own layer profile $\beta_\nu$, and
$\beta=\sum_{\nu\in\Sigma^\circ_b}\beta_\nu$. By the tail lemma
[c2, Lem.~4.1] (\texttt{tail-of-pf2-logconcave}, \texttt{proved}): \emph{if $\beta$ is
log-concave with interval support then $\gamma$ is log-concave}, and then by the radial
lemma [c2, Lem.~4.3] $G$ is $\PFtwo$. So condition~(A) at general $m$ follows from
\[
  (\mathrm B)\qquad \beta\ \text{is log-concave with interval support.}
\]
At $m=2$, [c2, Prop.~3.2] proves the much stronger statement that $\beta=c_+$ for an
explicit concave $c$. That is what dies.

\begin{proposition}[step 3 of {[c2]} fails at $m\ge3$]\label{prop:dead}
\begin{enumerate}
\item For a \emph{single} width vector $w\in\Z_{\ge1}^m$, the layer profile $\beta_w$ of
$*_i\mathbf 1_{[0,w_i-1]}$ is log-concave in all $37\,448$ cases with $m\le5$, $w_i\le8$,
but is \emph{not} the positive part of a concave function in $24\,192$ of them. The
smallest witness is $w=(3,3,3,3)$, where $\beta_w=(3,6,6,3,1)$ (innermost layer first):
log-concave, with increments $+3,0,-3,-2$, hence not concave.
\item For the \emph{aggregate} $\beta$ on a real slice, $\beta=c_+$ fails on $162$ of the
$44\,173$ nonzero slices with $m=3$, $n\le10$, $d\le12$. The smallest by $d$ is $n=9$,
$\mu=(0,1,6)$, $\lambda=(2,5,8)$, $b=2$, where $\Sigma^\circ_b$ has three elements:
\[
\begin{array}{c|c|c|c|c}
 \nu & (L_i,R_i)_i & w(\nu) & f_\nu & \beta_\nu\\\hline
 (0,1,8) & (0,0),(3,5),(8,8) & (1,3,1) & (1,1,1)\ \text{on }[2,4] & (0,1)\\
 (0,2,7) & (0,1),(3,5),(7,8) & (2,3,2) & (1,3,4,3,1)\ \text{on }[1,5] & (1,2,1)\\
 (0,3,6) & (0,2),(3,5),(6,8) & (3,3,3) & (1,3,6,7,6,3,1)\ \text{on }[0,6] & (1,3,2,1)
\end{array}
\]
all three concentric about $c=3$ (consistently with Theorem~\ref{thm:C}: $d=8$, $b=2$), with
$G=(1,4,10,12,10,4,1)$ and $\beta=(2,6,3,1)$. Each $\beta_\nu$ \emph{is} concave; the sum is
not, because the positive-part-of-concave class is not closed under addition when the
supports differ. $G$ is nevertheless $\PFtwo$ and $\beta$ is log-concave.
\item The pairwise route fails at the layer level too. On the instance of~(2),
$\beta_{(1,3,1)}=(0,1,0,0)$ and $\beta_{(3,3,3)}=(1,3,2,1)$ violate
$2g(r)h(r)\ge g(r-1)h(r+1)+g(r+1)h(r-1)$ at $r=2$: $0\ge1$. Over $m\le5$, $n\le9$,
$d\le11$ there are $2\,867$ failing layer pairs out of $84\,051$. So the antichain
obstruction of [c1, §2.1 item~1] is not removed by passing to layers.
\item $\beta$ is not real-rooted (e.g.\ $\beta=(1,1,1)$ at $n=6$, $\mu=(0,1)$,
$\lambda=(2,5)$, $b=2$), so no interlacing argument is available --- the same obstruction
as [c1, §2.1 item~3], one level down.
\end{enumerate}
\end{proposition}
\begin{proof}
All four are finite computations; see §\ref{sec:verif}. The structural remarks are checked
by hand on the displayed data: for (1), $(1,1,1)^{*4}=(1,4,10,16,19,16,10,4,1)$ has
$\gamma=(19,16,10,4,1)$ and $\beta=(3,6,6,3,1)$; for (2), $\gamma=(12,10,4,1)$ gives
$\beta=(2,6,3,1)$ with increments $+4,-3,-2$, and $(0,1)+(1,2,1)+(1,3,2,1)=(2,6,3,1)$.
\end{proof}

\section{The gap, stated so that it has no quantifier left to shrink}\label{sec:gap}

Theorems~\ref{thm:C} and~\ref{thm:tent} say that the slice family is a family of
\emph{concentric box splines whose half-widths are gap-free}. Does that abstract datum
suffice? No --- and the obstruction is not that the datum is too weak by a little. It is
that the datum is \emph{shape-blind}, and no shape-blind hypothesis can work.

\begin{proposition}[three successively stronger abstract hypotheses, all insufficient]\label{prop:abs}
Let $g_1,\dots,g_N$ be convolutions of interval indicators, all centred at $0$, with
half-widths $\Lambda_1,\dots,\Lambda_N$ and multiplicity profile $M(r)=\#\{j:\Lambda_j=r\}$.
Each $g_j$ is $\PFtwo$.
\begin{enumerate}
\item \textbf{Concentricity alone is insufficient.} $g_1=\delta_0$ $(w=(1))$ and
$g_2=\mathbf 1_{[-2,2]}$ $(w=(5))$ give $g_1+g_2=(1,1,2,1,1)$, not $\PFtwo$; here
$M=(1,0,1)$ has a gap.
\item \textbf{Adding gap-freeness is insufficient.} Three copies of $\delta_0$, one
$\mathbf 1_{[-1,1]}$ and one $\mathbf 1_{[-2,2]}$ give $(1,2,5,2,1)$, not $\PFtwo$, with
$M=(3,1,1)$: gap-free, but not log-concave, since $1^2<3\cdot1$.
\item \textbf{Adding ``$M$ is the positive part of a concave function'' is insufficient.}
At $m=2$ take $w^{(1)}=(1,1)$, $w^{(2)}=(2,2)$, $w^{(3)}=(1,5)$, $w^{(4)}=(2,6)$, i.e.
\[
  \delta_0,\quad (1,2,1),\quad (1,1,1,1,1),\quad (1,2,2,2,2,2,1),
\]
with half-widths $0,1,2,3$, so $M=(1,1,1,1)$ --- flat, hence concave. The sum is
$(1,3,4,6,4,3,1)$, and $4^2=16<3\cdot6=18$: not $\PFtwo$.
\end{enumerate}
\end{proposition}
\begin{proof}
Each item is the displayed arithmetic. (1) $\delta_0+\mathbf 1_{[-2,2]}=(1,1,2,1,1)$ and
$1^2<1\cdot2$ at the points adjacent to the centre. (2)
$3\delta_0+\mathbf 1_{[-1,1]}+\mathbf 1_{[-2,2]}=(1,2,5,2,1)$ and $2^2<1\cdot5$. (3) summing
the four displayed sequences after centring gives $(1,3,4,6,4,3,1)$.
\end{proof}

Item~(3) is not merely one more counterexample; it closes the whole family.

\begin{theorem}[no shape-blind hypothesis can suffice]\label{thm:blind}
There exist two multisets of concentric convolutions of interval indicators with the
\emph{same} half-width multiplicity profile $M=(1,1,1,1)$ --- the flattest and most concave
profile available --- such that one sum is $\PFtwo$ and the other is not:
\[
\begin{array}{lll}
  \{(1,1),(2,2),(1,5),(2,6)\} &\longrightarrow& (1,3,4,6,4,3,1)\ \ \text{\emph{not} }\PFtwo,\\[2pt]
  \{(1,1),(2,2),(3,3),(4,4)\} &\longrightarrow& (1,3,6,10,6,3,1)\ \ \PFtwo.
\end{array}
\]
Consequently no condition formulated in terms of the multiset $\{\Lambda_j\}$ alone --- no
matter how strong --- implies $\sum_jg_j\in\PFtwo$. Any proof of condition~(A) at $m\ge3$
must use the \emph{width vectors} $w(\nu)$ themselves, not only the half-widths
$\Lambda(\nu)=\tfrac12\bigl(\sum_iw_i(\nu)-m\bigr)$ they determine.
\end{theorem}
\begin{proof}
Both multisets have half-widths $0,1,2,3$, hence the same $M$. The first sums to
$\delta_0+(1,2,1)+(1,1,1,1,1)+(1,2,2,2,2,2,1)=(1,3,4,6,4,3,1)$, which fails $\PFtwo$ at
$16<18$. The second sums to
$\delta_0+(1,2,1)+(1,2,3,2,1)+(1,2,3,4,3,2,1)=(1,3,6,10,6,3,1)$, which satisfies
$6^2=36\ge3\cdot10=30$ and $10^2\ge6^2$, and has interval support. A hypothesis that is a
function of $M$ alone takes the same value on both, so it either admits the first (and is
not sufficient) or excludes the second (and is not satisfied by all real slices, since
$M=(1,1,1,1)$ does occur).
\end{proof}

\begin{remark}[what the width vectors supply at $m=2$, and the shape of the open question]
Theorem~\ref{thm:blind} explains why the $m=2$ proof could not be shape-blind and was not.
There the width pair obeys $w_1(\nu)-w_2(\nu)=2(t_0-\nu_1)$ ([c2, eq.~(6)]), so the
\emph{difference} determines the slice parameter: no two summands of one slice can share it.
The first multiset of Theorem~\ref{thm:blind} has differences $0,0,-4,-4$ and is therefore
not realisable at $m=2$ --- the coupling between shape and position is exactly what rules it
out. The open question at $m\ge3$ is thus not ``which condition on $\Lambda$'' but:
\begin{quote}
  \emph{which $m$-tuples $w(\nu)$ can co-occur in one slice, and why does that family
  sum to a $\PFtwo$ sequence?}
\end{quote}
Each $w_i(\nu)=\min(\nu_{i+1}-1,\lambda_i)-\max(\nu_i,\lambda_{i-1}+1)+1$ is concave in each
of the two coordinates it sees, and $\sum_iw_i$ is pinned by Theorem~\ref{thm:C}; what is
missing is a usable description of the joint law of the $w_i$ along a slice.
\end{remark}

\begin{remark}[the route through $(\mathrm B)$ is strictly too strong]\label{rem:toostrong}
It is tempting to prove $(\mathrm B)$ --- $\beta$ log-concave --- and quote the tail lemma,
and $(\mathrm B)$ does hold on every real slice tested (§\ref{sec:verif}): $4\,323\,908$
slices, including $m=3$ with $n\le16$, $d\le20$. But $(\mathrm B)$ is \emph{false} in the
abstract class of Proposition~\ref{prop:abs}(2) while the target is \emph{true} there: for
$N=2$, $m=2$, $w^{(1)}=(1,5)$ and $w^{(2)}=(3,3)$ (half-widths $2,2$, gap-free), the sum is
$(1,1,1,1,1)+(1,2,3,2,1)=(2,3,4,3,2)$, which \emph{is} $\PFtwo$, while $\beta=(1,1,2)$ is
\emph{not} log-concave, since $1^2<1\cdot2$. In the exhaustive abstract search of
§\ref{sec:verif}, $(\mathrm B)$ fails on $23$ gap-free multisets at $m=2$ and $198$ at
$m=3$, and the $\PFtwo$ conclusion fails on \emph{none} of those. So $(\mathrm B)$ cannot be
derived from Theorems~\ref{thm:C} and~\ref{thm:tent}: it is a true statement about real
slices that is strictly stronger than what those two theorems carve out, and proving it
requires the same missing input as proving (A) directly, plus more. For a proof it is
therefore the wrong intermediate; for a \emph{test} it is a good one, because it is sharper
than (A) and so falsifies faster.
\end{remark}

\begin{remark}[where the concavity went]\label{rem:Mconcave}
One positive structural fact survives and is worth recording even though
Theorem~\ref{thm:blind} shows it cannot be sufficient on its own: on real slices the
half-width multiplicity profile $M$ \emph{is} the positive part of a concave function, on all
$135\,422$ nonzero slices with $2\le m\le5$, $m\le n\le m+5$, $d\le11$. So the concavity that
the $m=2$ proof spends at the \emph{layer} level --- where it is false at $m\ge3$,
Proposition~\ref{prop:dead} --- reappears one level up, at the half-width level, where it is
a consequence to be proved from Theorem~\ref{thm:tent} rather than an input. By
Theorem~\ref{thm:blind} it will have to be combined with shape information, but it is a
genuine invariant of the family and the first thing I would try to prove next.
\end{remark}

\section{The control}\label{sec:control}

The brief asked for a refusal control: break the cylindric alignment and confirm that (C)
fails. The proof of Theorem~\ref{thm:C} says more than that --- it says \emph{which}
deformations must break (C) and which must not --- so the control is run as a four-way
panel with all four outcomes predicted in advance from \eqref{eq:tele}. Deform
\eqref{eq:setup} to
\[
  L_i^{(p)}(\nu)=\max(\nu_i,\ \lambda_{i-1}+p),\qquad
  R_i^{(q)}(\nu)=\min(\nu_{i+1}-q,\ \lambda_i),
\]
and let the $\lambda$-period be $n'$ while the $\nu$-period stays $n$.

\begin{center}
\begin{tabular}{llll}
 & deformation & prediction & $m=3$ outcome\\\hline
P1 & $p=q=2$ and $p=q=4$ (offset \emph{value}) & silent & \textbf{(C) holds}\\
P2 & $p=2,q=1$ and $p=1,q=3$ (offset \emph{mismatch}) & fires & fails $792/1306$, $765/1306$\\
P3 & $n'=n\pm1$, $n'=n+2$ (periods decoupled) & fires & fails $295/1306$, $396$, $377$\\
P4 & no wrap: $L_1=\nu_1$, $R_m=\lambda_m$ (ordinary skew) & silent & \textbf{(C) holds}\\
\end{tabular}
\end{center}

All four predictions hold, at $m=2,3,4$ (counts in §\ref{sec:verif}). P1 is silent because
Lemma~\ref{lem:adj} holds for every matched offset; P2 fires because the local collapse is
destroyed; P3 fires because the unmatched boundary term $R_m-R_0$ becomes $\nu$-dependent;
P4 is silent because in the linear case the two boundary terms degenerate \emph{separately}
--- $L_1=\nu_1$ is absorbed into $S_\nu$ and $R_m=\lambda_m$ is a constant --- which is the
statement that ordinary skew Kostka numbers also have concentric slice sums.

A control that fires on everything would not have distinguished the offset match from the
periodicity, and it was P3 that corrected my first (wrong) attribution; a control that
fired on nothing would have said the hypotheses are not load-bearing. This one fires on
exactly the two deformations the proof consumes.

\section{Verification}\label{sec:verif}

Code in \verb|proofs/code-1002-m3/|. No instrument's output was believed before it
reproduced a value already known.

\paragraph{Calibration.}
\begin{itemize}
\item The general-$m$ slice-sum engine reproduces both known $m=2$ values exactly, with
offsets: $k(\cdot,2)=(1,3,6,7,6,3,1)$ at $n=8$, $\mu=(0,3)$, $\lambda=(4,7)$
([c1, §``Concavity is false'']) and $G=(1,4,7,10,11,10,7,4,1)$ at $n=10$, $\mu=(0,5)$,
$\lambda=(5,12)$, $b=4$ ([c2, Prop.~2.4]).
\item \textbf{Estimand check.} An independent enumerator builds the chains
$\mu\prec\nu\prec\kappa\prec\lambda$ directly from \verb|is_shape|, \verb|hstrip|,
\verb|size| of \verb|proofs/code-q254/cyl.py|, never touching $L_i,R_i$. It agrees with
\eqref{eq:slice} on all $13\,689$ slices tested: $m=2$ ($n\le7$, $d\le9$; $2415$), $m=3$
($n\le7$, $d\le8$; $5124$), $m=4$ ($n\le7$, $d\le7$; $4150$), zero disagreements. The same
independent table confirms $k(a,b)=k(b,a)$ \emph{and} $k(a,b)=k(d-a-b,b)$ on $14\,833$
entries, zero bad --- the latter being the total-sum shadow of Theorem~\ref{thm:C}.
\item \textbf{Refusal probes.} The $\PFtwo$ test refuses $(1,1,2,1,1)$; the
concave-positive-part test refuses $(3,6,6,3,1)$; the four perturbed forms of
\eqref{eq:Lam} and \eqref{eq:effbox} ($\lambda_{i-1}+1\mapsto\lambda_{i-1}+0,+2$;
$\lambda_{i-2}+2\mapsto\lambda_{i-2}+1,+3$) are refused on $229/466$, $466/466$,
$400/1047$, $434/1047$ instances respectively.
\end{itemize}

\paragraph{Results.}
\begin{itemize}
\item \textbf{Theorem~\ref{thm:C}, STEP 0 census.} $\sum_i(L_i+R_i)$ is constant on every
one of $8\,335\,376$ nonempty slices, and equals $S_\nu+\abs{\lambda}$ on every one of
$40\,735\,322$ summands --- the closed form, not merely constancy --- over
$m=2$ ($n\le16$, $d\le20$), $m=3$ ($n\le14$, $d\le18$), $m=4$ ($n\le14$, $d\le18$),
$m=5$ ($n\le12$, $d\le14$), $m=6$ ($n\le11$, $d\le12$). The range is deliberately wider
than the $m=2$ range of [c2] in every parameter.
\item \textbf{Theorem~\ref{thm:tent}.} \eqref{eq:Lam} is exact on all $43\,916$ nonzero
$f_\nu$ and \eqref{eq:effbox} is exact on all $94\,906$ $\nu$ tested ($2\le m\le5$,
$n\le m+4$, $d\le8$). Gap-freeness of $\Lambda(\Sigma^\circ_b)$: $78\,312$ of $78\,312$
slices ($2\le m\le5$, $n\le9$, $d\le11$).
\item \textbf{Proposition~\ref{prop:dead}.} (1) $37\,448$ width vectors, $m\le5$,
$w_i\le8$: $\beta_w$ log-concave in all, $=c_+$ in $13\,256$ and not in $24\,192$. (2)
$44\,173$ slices at $m=3$, $n\le10$, $d\le12$: $\beta=c_+$ fails on $162$. (3) $2\,867$
failing layer pairs of $84\,051$. (4) $\beta$ real-rooted on $74\,669$ of $78\,312$.
\item \textbf{Condition (A) and (B), range-extended.} $(\mathrm A)$ and $(\mathrm B)$ both
hold with zero failures on every slice tested: $m=3$, $n\le16$, $d\le20$
($1\,519\,862$ slices); $m=4$, $n\le13$, $d\le16$ ($1\,502\,376$); $m=5$, $n\le12$,
$d\le13$ ($1\,069\,479$); $m=6$, $n\le11$, $d\le12$ ($232\,191$); $m=7$ in progress, see
\verb|code-1002-m3/wide.out| for the running counts. ($\beta=c_+$ for a concave $c$ fails
on $105\,084$ of the $m=3$ slices, $57\,877$ of the $m=4$, $14\,870$ of the $m=5$ and $60$
of the $m=6$ ---
Proposition~\ref{prop:dead}(2) at scale.) This extends the registry's
$96\,781$ interior points ($n\le9$, $m\le7$, $d\le11$) by more than an order of magnitude
in $m=3,4$. $G$ was symmetric and $\beta\ge0$ on every slice, as
Theorem~\ref{thm:C} requires.
\item \textbf{Proposition~\ref{prop:abs} and Remark~\ref{rem:toostrong}.} Exhaustive over
multisets of concentric box splines: at $m=1$ with $w\le9$ and size $\le6$, gap-freeness
gives zero $\PFtwo$ failures at sizes $2,3,4$ and \textbf{13 failures at size 5}, \textbf{32
at size 6}; at $m=2$ with $w_i\le5$, zero at sizes $2,3$ and $5$ failures at size $4$, $21$
at size $5$. \emph{The discovery set was sizes $\le3$, where the score was perfect.} The
$(\mathrm B)$ failures in the gap-free class are $23$ at $m=2$ and $198$ at $m=3$, with the
$\PFtwo$ conclusion failing on none of those.
\item \textbf{Theorem~\ref{thm:blind} and Proposition~\ref{prop:abs}(3).} Restricting the
abstract search to multisets whose $M$ is the positive part of a concave function: $m=1$,
$w\le11$, sizes $2$--$8$ ($20,36,50,71,104,127,167$ multisets) gives \textbf{zero} $\PFtwo$
failures; $m=2$, $w_i\le6$, sizes $2,3$ ($216$, $1238$) gives zero, and size $4$
($4606$) gives \textbf{4 failures}, all equivalent to
Proposition~\ref{prop:abs}(3) under the symmetry $w\mapsto$ reversal; $m=3$, $w_i\le5$,
sizes $2,3$ ($2444$, $45\,369$) gives zero. Again the discovery set --- sizes $\le3$ ---
scores perfectly and the statement is false at size $4$. The two multisets of
Theorem~\ref{thm:blind} were then checked directly.
\item \textbf{Proposition~\ref{prop:abs}(3).} $M=c_+$ for a concave $c$ on all $135\,422$
nonzero slices with $2\le m\le5$, $m\le n\le m+5$, $d\le11$; $M$ fails to be nondecreasing on $10\,000$ of them and fails to be
nonincreasing on $11\,966$, so the concavity is not an artefact of monotonicity.
\end{itemize}

\section{Gaps, and what comes next}

\begin{enumerate}
\item \textbf{No gap in Theorem~\ref{thm:C} or Theorem~\ref{thm:tent}.} Both are proved
from Lemma~\ref{lem:adj}, the cyclic conventions, and the connectivity of a box-slice. The
one external input is [c1, Prop.~2.3], \texttt{proved}.
\item \textbf{Condition~(A) at $m\ge3$ is still open, and is now known not to be
shape-blind.} Theorem~\ref{thm:blind} rules out every hypothesis phrased in the half-widths
alone, so Theorems~\ref{thm:C} and~\ref{thm:tent} are provably not a complete reduction.
Two concrete next targets, in order: (i) prove that $M$ is the positive part of a concave
function (Remark~\ref{rem:Mconcave}) --- $M(r)$ counts the lattice points of a box-slice on
which the separably convex $\sum_i\phi_i(\nu_i)$ takes a fixed value, and $\Lambda$ is
$1$-Lipschitz along exchanges by Theorem~\ref{thm:tent}(3), so this should be accessible;
(ii) describe which width vectors $w(\nu)$ co-occur on a slice, which is the input
Theorem~\ref{thm:blind} says is unavoidable.
\item \textbf{Three steps of the $m=2$ proof, at $m\ge3$.} Step~1 (concentricity):
\textbf{proved}, Theorem~\ref{thm:C}. Step~2 (the tent): \textbf{proved},
Theorem~\ref{thm:tent}. Step~3 ($\beta=c_+$ for concave $c$): \textbf{false},
Proposition~\ref{prop:dead}, and the concavity relocates to $M$ --- where, by
Theorem~\ref{thm:blind}, it is still not sufficient. Step~4 (tail lemma + radial lemma):
\textbf{unchanged}, both are statements about one sequence and are already proved for all
$m$; they apply the moment $(\mathrm B)$ is available, but
Remark~\ref{rem:toostrong} shows $(\mathrm B)$ is the wrong thing to aim at.
\item \textbf{Condition (B)} (\texttt{conj-B-logsubmodular}) is untouched and remains
\texttt{computed}.
\item \textbf{A by-product worth recording.} Control P4 shows that concentricity holds for
\emph{ordinary} skew shapes too, by a degeneration of the same telescoping. So the
mechanism is not cylindric; what is cylindric is only the form of the boundary term.
\end{enumerate}

\section{Summary}

The brief asked whether concentricity is an $m=2$ accident or a cylindric fact. It is
neither: it is a \emph{horizontal-strip} fact. One scalar identity,
$\max(x,y+1)+\min(x-1,y)=x+y$, applied once per bead, plus one global closure of the cycle,
gives it at every $m$, and a four-way control confirms that attribution by firing on exactly
the two deformations the proof consumes. The tent lemma generalises too, and more cleanly
than at $m=2$: the half-width is a separably concave function on a box-slice, so it is
gap-free for the standard exchange-connectivity reason.

What does not generalise is the certificate, and the session's sharpest finding is
negative. The layer profile is no longer the positive part of a concave function; the route
through its log-concavity is strictly stronger than the target and false in the abstract
class the two theorems carve out; and the concavity reappears one level up, at the
half-width multiplicity profile --- where Theorem~\ref{thm:blind} shows it \emph{still}
cannot suffice, because two multisets with the same half-width profile can land on opposite
sides of $\PFtwo$. That last result is the one I would keep: it says the two theorems proved
here are provably not a complete reduction, and that any proof of condition~(A) at $m\ge3$
must know which width vectors co-occur in a slice. I had intended to leave a conjecture in
its place; the conjecture was refuted by the search I launched to test it, at multiset size
$4$ after sizes $2$ and $3$ scored perfectly. That is the third session in a row in which
extending a range past the one that bore the statement was the decisive move, and the first
in which it caught a conjecture of mine before I wrote it down as one.

\end{document}
